%
%============= CHAPTER BREAK =================%
%
\chapter{Herodian Gramm. (ca 161--180 CE)}%
\vspace*{3mm}%
\justifying
%
\noindent{}Herodian the Grammarian should not be confused with the early 3rd century CE Herodian the Historian. %
The former was the son and fellow grammarian of Apollonius Dyscolus. %
Scholars accept the authenticity of several works by this Herodian, %
chief of which is the \textit{General Prosody} / \textit{De prosodia catholica}, %
with some caveats as to the reliability of the most recent scholarly edition (Lentz 1867).\endnote{ %
For an example of using Herodianʼs \textit{General Prosody} %
as an historical reliable source for onomastics, %
see Bremmer 1998:156, who finds the proper name Aëtius stemming originally from Pausanias, %
yet still trustworthy as an attestation in Herodianʼs \textit{Pros.Cath.} 1.119.35--120.6, %
a conclusion with which Mansfeld 2016:166 concurred.% 
} % %
By comparison, the \textit{Word Divisions} / \textit{Partitiones} / Ἐπιμερισμοί traditionally attributed to Herodian are, as Dyck put it, %
“surrounded by misunderstandings”. %

Foremost among current misunderstandings is whether the \textit{Partitiones} attributed to Herodian %
should be considered as a single work in two books, or as two separate works. %
Assembling a comprehensive catalog of works attributed to Herodian (correctly or falsely), %
Dickey cataloged these as two distinct works called the \textit{Partitiones} / %
Ἐπιμερισμοί / \textit{Epimerismoi} that traveled under the name of Herodian.\endnote{%
Dickey 2014:329.} %
Dickey claimed that the first of these, not available in the TLG, %
has only fragmentary witnesses by later authors, %
may originally have derived from Herodian,
but has been so suffused with later glosses as to make the recovery of %
any authentic Herodian content impossible and the question of authenticity moot. %
Dickey cataloged the second \textit{Partitiones} / \textit{Epimerismoi} as a lengthy and intact work %
available in a scholarly (however dated) edition,\endnote{%
Boissonade 1819, based only on codd. Paris 2543 and 2570. Pinakes currently lists 21 extant manuscripts.} %
one available to search in the TLG, %
yet entirely spurious in authorship. %

Lehrs (1848) and Lentz (1867--1870) had previously found the long-form \textit{Partitiones} spurious, %
mainly based on the title of the work as indicating a setting of school instruction, %
unbefitting the serious work of a renowned grammarian.\endnote{%
Dyck 1981:225 noted the weakness of this argument.} %
Yet Lentz also conceded that the doctrines of the longer composition were authentically Herodian %
and likely contained some authentically Herodian source material, even if thoroughly excerpted and epitomized. %
Dyck demonstrated the reception of fragments of the bivolume \textit{Partitiones} %
in the 10th century CE \textit{Epimerismi Homerici}\endnote{Dyck 1981.} %
as well as in the 6th century CE \textit{Etymologicum} of Orion of Thebes.\endnote{Dyck 1993:792.} %
In Dyckʼs opinion, at least one part of the two-part work %
had an alphabetical order, %
but both parts are so saturated with glosses from later grammarians, %
such as Orus, Helladius, John Philoponus, and Zenobius, %
that “any attempt to find a Herodianic ‘core’ is doomed”.\endnote{%
Dyck 1993:793.}
In his \textit{General Prosody}, however, %
Herodian claims to have indexed some 60,000 different Greek words.\endnote{%
Forbes et al 2016.} %
It seems to me an open question whether the original substrate of the long-form \textit{Partitiones}, %
particularly given its simplicity in concept and detail, %
originally comprised just such a word index, %
suitable for reference and perhaps beginner levels of education. %

NT scholarship does not often consider the \textit{Partitiones} (long-form or fragmentary) as potentially significant. %
Even when this does happen, %
sometimes the approach is too naïve about authenticity or too cut and dry about inauthenticity. %
Let us take one excerpt as an example. %

\vspace*{3mm} 
 \begin{quote}Πᾶσα λέξις ἀπὸ τῆς βη συλλαβῆς ἀρχομένη διὰ τοῦ η γράφεται· %
 οἷον· Βηθλεὲμ, καὶ Βηθανία, Βηθεσδὰ, Βηθσαϊδὰ, %
 καὶ Βηθσφαγὴ, ὀνόματα τόπων ἐν Ἱεροσολύμοις· Βὴλ, θεὸς Αἰγύπτιος· (Boissonade 1819:5)\\[1mm] %
 Every word that begins with the syllable βη is written with the η, %
 such as: Bethlehem, and Bethany, Bethesda, Bethsaida, %
 and Bethphage, names of places in Jerusalem. Bel, an Egyptian God. %
 \end{quote} % 
 \vspace*{3mm}
 
\noindent{}In regard to the former, Smith takes this excerpt from the long-form \textit{Partitiones} as authentic,\endnote{%
Smith 2004.} %
when it likely reflects later Christian interpolation %
and Christian interest in onomastics of the sort first seen in Eusebius of Caesarea. %
On the other hand, though Kučera is almost certainly correct that the excerpt is inauthentic, %
the dismissal of the entire long-form \textit{Partitiones} as inauthentic is over-simplistic overreach.\endnote{%
Kučera 2025:268, citing to Dickey 2014 and Dyck 1981.}

While I grant that later glosses cloud both books of the \textit{Partitiones}, %
and that reconstruction of an authentic core for either is not achievable at present (nor within the scope of this work), %
it is prejudicial to rule out any and all potentially authentic Herodian content in the long-form \textit{Partitiones}. %
That is to say, for the purposes of establishing the emergence of a neologism, %
each excerpt should be evaluated on its own merits, %
compared with Herodianʼs clearly authentic writings, %
and also set within the broader diachronic sweep of a lemmaʼs history and usage. %

Thus if Herodianʼs father, Dyscolus, is the first in extant, digitized Greek literature to use a certain term, %
\textit{and} the same term is found in the long-form \textit{Partitiones}, %
\textit{and} the hand of a later gloss is not clearly in evidence, %
then it is unreasonable to posit such an excerpt as a later interpolation. %
Indeed, the question of authenticity vs. inauthenticity could run in reverse, %
since later glossers, epitomizers, and compilers routinely mixed up excerpts of Dyscolus and Herodian, %
which is to say that a given Herodian attribution may deserve consideration as a Dyscolan excerpt. %

When a Greek neologism appears first in the authentic \textit{General Prosody}, %
and also appears in the long-form \textit{Partitiones}, %
then again authenticity is reasonable, unless the hand of a later gloss is demonstrable. %
To go a stretch further, if the long-form \textit{Partitiones} contains what may be the first instance of a Greek word, %
and late 2nd century non-Christian authors (i.e., rough contemporaries of Herodian) also attest to this word, %
then it is reasonable to hold that the word reflects an authentic neologism that emerged around the time of Marcus Aurelius. %
That is to say, in this book we are primarily interested in \textit{chronological-attributional precision} %
more than \textit{authorial authenticity/disambiguation}. %

However, if a word is first attested in non-Christian writings only in the long-form \textit{Partitiones}, %
and not by any other non-Christian writer in the late 2nd or early 3rd century, %
then there is reason to doubt and leave the attribution ambiguous as [Ps-]Herodian. %
Be that as it may, when and if examples reflect later interpolations, %
then that only pushes the likely emergence of such neologisms later in history, %
i.e., after the time of Marcus Aurelius rather than contemporaneous to him. %
Herodianine authenticity strengthens the overall thesis of the book that the NT %
contains a variety of mid-2nd to early-3rd century Greek neologisms, %
but Herodianine inauthenticity strengthens it even more. %
The diachronic approach to pinpointing neologisms may even model a new, %
fruitful line of research in regard to delineating between more and less probable Herodianine expressions. \\[1.5mm] %

\noindent{}\textbf{ἀνάγαιον} / \textit{anagaion} / “above-ground” derives from the classical lemma %
ἀνώγεον / \textit{anōgeon} / “above-ground”.\endnote{%
The earlier form is used by Attic Xenophon (\textit{Anab}. 5.4.29), Antiphanes (fr. 312), et al.} %
Philo uses a spelling variant, %
ἀνάγειον / \textit{anageion}, and %
elaborates on the termʼs architectural significance; %
he specifically notes that women build these raised areas %
between houses with surrounding walls but no roofs, %
so as to guard their modesty and allow their voices to travel upwards unimpeded.\endnote{%
\textit{Vit.contempl.} 33.} %

The first non-Christian use of the spelling ἀνάγαιον / \textit{anagaion} %
may have been found in the long-form \textit{Partitiones} of [Ps-]Herodian.\endnote{%
Boissonade 1819:209. Note that several other forms also appear in Herodianʼs %
\textit{De orthographia} (Lentz 1870), such as ἀνόκαιον (p. 476), ἀνώγαιον (p. 596), %
ἀνώγειος (p. 440, 477), and ἀνώγεων (p. 440, 476).} %

\vspace*{2mm} 
 \begin{quote}Διὰ διφθόγγου δὲ καὶ ο μικροῦ ταῦτα· %
 ὑπόγαιος· \textbf{ἀνάγαιος}· κατάγαιος. \\[1mm] %
 But these [are constructed] with the diphthong and the short o: %
 upon-ground; \textbf{above-ground}; below-ground. \end{quote} % 
\vspace*{2mm}

\noindent{}The so-called upper-room setting where Jesus is said to have shared the Passover %
 with his students is described with this same term in canonical Mark and canonical Luke. %
 The latter appears in a group of verses (22:9--13) unattested and very likely missing from Marcion's \textit{Evangelion}. %
 Both canonical texts follows the apparently new and unique spelling of Herodian. %
 While there are several spelling variants in the manuscripts of Mark and Luke, %
 this particular spelling clearly has the earliest and strongest manuscript attestation:\endnote{%
 E.g., ἀνώγεον is found in Γ Δ Ω 2 157 \textit{ƒ}\textsuperscript{1} of Mark and Γ 1 118 of Luke.} %

\vspace*{2mm} 
 \begin{quote}\textls[-15]{καὶ αὐτὸς ὑμῖν δείξει \textbf{ἀνάγαιον} μέγα ἐστρωμένον ἕτοιμον·} (Mark 14:15) \\ %
 And he will show you a large \textbf{above-ground} level prepared. \end{quote} % 
\vspace*{1mm} 
 \begin{quote}κἀκεῖνος ὑμῖν δείξει \textbf{ἀνάγαιον} μέγα ἐστρωμένον· (Lk 22:12) \\ %
 Eʼen he will show you a large \textbf{above-ground} level. \end{quote} % 
 \vspace*{2mm}

Based on this lexicographical-chronological finding, %
the portraits of Jesus piously sharing a traditional ritual meal with his disciples, %
one that foreshadowed and guaranteed his innocent non-violence, %
is best read as a literary product of the time of Marcus Aurelius or later.\endnote{%
The violence of the young Jesus in \textit{InfThom} recension 3, 9.1 \textit{bis} is interesting to compare here, %
since it is from a level “above-ground” / ἀνωγαίου|ἀνώγαιον that a childhood rival is killed by Jesus.} \\[1mm] %

\noindent{}\textbf{ἀπροσωπόλη[μ]πτος} / \textit{aprosōpolē[m]ptos} / “unbiased by person” (adj), \\
\noindent{}\textbf{ἀπροσωπολή[μ]πτως} / \textit{aprosōpolē[m]ptōs} / “unbiased by person” (adv), %
including various derivative nouns and verbs, %
first seem to appear in non-Christian writings in the long-form \textit{Partitiones} of [Ps-]Herodian.\endnote{%
Boissonade 1819:188, 269.} %

\vspace*{2mm}
\begin{quote}\textls[-15]{Εὔεικτος δὲ γράφεται διὰ τοῦ εἰ διφθόγγου. %
Τὰ παρὰ τὸ λήβω, τὸ λαμβάνω, γράφονται διὰ τοῦ η· %
οἷον: ἄληπτος, ἀκράτητος, ἀκατάληπτος, ἀνεπίληπτος, %
ἐπίληπτος, \textbf{ἀπροσωπόληπτος}, φρενόληπτος·} \\[1mm] %
Now “pliant” (\textit{eueiktos}) is written with the diphthong εἰ. %
Those [derived] from “to take” and “to receive” are written with η, %
such as: ungraspable, uncontrollable, incomprehensible, unassailable, % 
culpable, \textbf{unbiased}, mentally-composed.
\end{quote}
\clearpage

\begin{quote}
Ἄληπτος, ἀκατάληπτος, ἀπερίληπτος, \textbf{ἀπροσωπόληπτος}, %
ἀκατα\-ληψία, ἐπιληψία, καὶ γίνεται ἐκ τοῦ λήβω. \\[1mm] %
Ungraspable, incomprehensible, uncircumscribed, \textbf{unbiased}, %
incom\-prehension, stoppage, are also derived from “to take”. %
\end{quote}
\vspace*{2mm}

\noindent{}While [Ps-]Herodian only uses an adjectival form, %
the NT only contains an adverbial form, and one instance of the same, %
at 1Pet 1:17: %

\vspace*{3mm} 
 \begin{quote}\textls[-15]{καὶ εἰ πατέρα ἐπικαλεῖσθε τὸν \textbf{ἀπροσωπολήμπτως} κρίνοντα κατὰ τὸ ἑκάστου ἔργον, %
 ἐν φόβῳ τὸν τῆς παροικίας ὑμῶν χρόνον ἀναστράφητε.} \\[1mm] %
 And if you invoke as Father him who judges \textbf{without bias}, based on the work of each, %
 conduct the time of the sojourn of yours in fear. \end{quote} % 
 \vspace*{3mm}

\noindent{}A similar, adverbial refrain (without the medial μ) appears in \textit{EpBarn} 4:12: %

\vspace*{3mm} 
 \begin{quote}Ὁ κύριος \textbf{ἀπροσωπολήπτως} κρινεῖ τὸν κόσμον· %
 ἕκαστος καθὼς ἐποίησεν κομιεῖται—ἐὰν ᾖ ἀγαθός, %
 ἡ δικαιοσύνη αὐτοῦ προηγήσεται αὐτοῦ· %
 ἐὰν ᾖ πονηρός, ὁ μισθὸς τῆς πονηρίας ἔμπροσθεν αὐτοῦ. \\[1mm] %
 The Lord will judge the world \textbf{without bias}. %
 Each will receive just as one did --- if one be good, %
 that oneʼs justice will precede oneself. %
 If one be bad, the reward of evil [will go] before oneself. \end{quote} % 
 \vspace*{3mm}

Though the emphasis on divine judgment without personal bias %
is distinctive to the early Christian examples, %
these still likely date during or after the reign of Marcus Aurelius, %
when the shared stem apparently first gained traction. \\[2mm] %

\noindent{}\textbf{ἀρτέμων} / \textit{artemōn} / “foresail” in its root form is spelled identically %
to the named entity Ἀρτέμων / \textit{Artemōn} / “Artemon”, %
only differing by way of the capitalization of the first letter, %
an optional custom among postclassical Greek scribes. %
The named entity may go back as far as Protagoras and Xanthus in the 5th century BCE. %
It is the name of a known Greek historian in the 2nd century BCE. %
Prior to the 1st century CE, this named entity appears 14 times in the TLG and 25 times in digitized Greek papyri, %
and some 10 additional times in Greek papyri in or after the 1st century CE.\endnote{%
In the papyri.info results view, P.Herc. 1044 does not quote the specific form, %
but it is Ἀρ-\/τέμ̣[ω]νι (fr. 33, lines 6--7).} %
It appears once in a Judean ossuary (JERU0139) inscribed between 20 BCE and 70 CE.\endnote{%
“Of Gaius, of Artemon, Berenikaian” / Γαίου Ἀρτέμωνος Βερνικέος; %
see \url{search.inscriptionsisraelpalestine.org/inscriptions/jeru0139}.} %
In the TLG, it appears once in Plutarch, writing in the late 1st or early 2nd CE, %
twice in a possibly Aurelian-era excerpt from the long-form \textit{Partitiones} of [Ps-]Herodian, %
and six times in the \textit{Deipnosophists} of Athenaeus in the early 3rd century CE. %
It even appears once in the NT, in Titus 3:12, %
as part of a litany of other likely fictitious names evoking learned allusions: %

\vspace*{3mm} 
 \begin{quote}Ὅταν πέμψω Ἀρτεμᾶν πρὸς σὲ ἢ Τύχικον, %
  σπούδασον ἐλθεῖν πρός με εἰς Νικόπολιν, %
  ἐκεῖ γὰρ κέκρικα παραχειμάσαι. \\[1mm] %
 Whenever I send Artemas unto you, or Tychicus, %
 hurry to come unto me in Nicopolis, %
 for there I have decided to winter. %
 \end{quote} % 
\vspace*{3mm}

The first non-Christian use of the distinct lemma ἀρτέμων / \textit{artemōn} / “foresail” as a common noun, %
not a proper noun, %
appears in the long-form \textit{Partitiones} of [Ps-]Herodian. %
After listing words that end in αινος / \textit{ainos}, then αιμων / \textit{aimōn}, %
all examples written with the diphthong αι / \textit{ai}, %
he then begins to note exceptions, starting with a list of words ending in εμων / \textit{emōn}: %

\vspace*{3mm} 
 \begin{quote}Πλὴν τοῦ \textbf{ἀρτέμων}· ἀκρέμων· Πολέμων· καὶ ἡγεμὼν,
 καὶ κηδεμών· ὀξύνονται. \\[1mm] %
 Excepting: \textbf{main pulley}|\textbf{foresail} (\textbf{\textit{artemōn}}); hanging (\textit{akremōn}); %
 Polemon; both hegemon and guardian (\textit{kēdemōn}) are acute-accented. %
 \end{quote} % 
\clearpage

The next non-Christian author to use the common noun is apparently the 5th or 6th century %
Hesychius the Alexandrian, in his roughly 50,000 word \textit{Lexicon}.\endnote{%
On the date of the work, see Dickey 2007:88.} %
The LSJ entry for the Greek common noun mentions its derivation from the Latin word %
\textit{artemon}, used by the mechanical engineer Vitruvius (1st cent. BCE) in his work \textit{De architectura} 10.2.9 %
to describe a key component of a complex crane.\endnote{%
OLD gives the definitions: %
1. The main block of a tackle. %
2. A jib or foresail. %
LSJ, s.v. ἐπάγων, cites to Heschychius the Alexandrian lexicographer, %
who equates this term to ηἀφνήδιος / \textit{ēafnēdios} and κασσάνδρα / \textit{kassandra}, %
both of which are apparently Byzantine, rather than ancient or late ancient terms.%
} %

\vspace*{2mm} 
 \begin{quote}In radice autem machinae conlocatur tertia troclea; %
 eam autem Graeci epagonta, nostri \textbf{artemonem} appellant. \\[1mm] %
 Yet at the base of the machine a third pulley-block is affixed; %
 the Greeks call it \textit{epagonta}, but we [call it] \textbf{\textit{artemon}}. \end{quote} % 
 \vspace*{2mm}

The common noun lemma ἀρτέμων / \textit{artemōn} / “foresail” %
is an NT \textit{hapax}, appearing only in the sea voyage of Paul in Acts 27:40: %

\vspace*{2mm} 
 \begin{quote}καὶ τὰς ἀγκύρας περιελόντες εἴων εἰς τὴν θάλασσαν, %
 ἅμα ἀνέντες τὰς ζευκτηρίας τῶν πηδαλίων, %
 καὶ ἐπάραντες τὸν \textbf{ἀρτέμωνα} τῇ πνεούσῃ κατεῖχον εἰς τὸν αἰγιαλόν. \\[1mm] %
 And after discarding the anchors, allowing them [to sink] into the sea, %
 simultaneously releasing the yoke-fittings of the rudders, %
 and hoisting the \textbf{foresail} to the wind, %
 they held tight for the shore.\end{quote} % 
 \vspace*{2mm} 

\noindent{}Incidentally, the lemma ζευκτήριος / \textit{dzeukērios} / “yoke-fit”, %
prior to its single appearance in the NT, %
only appears in the writings of Aeschylus, %
and does so at least three different times.\endnote{
\textit{Fragmenta}, Tetralogy 44 play A, fr. 736a // fr. 382 %
addressing “Father Wine-God” / πάτερ Θέοινε as “yoke-fitter of the frenzied” / μαινάδων ζευκτήριε. %
\textit{Persae} 736, where a Messenger tells Darius that Xerxes was “glad to cross the bridge, %
one a dual yoke-fitting” / ἄσμενον μολεῖν γέφυραν, ἓν δυοῖν ζευκτήριον. %
\textit{Agam} 529, describing Agamemnon as: % 
 \begin{quote}he who a yoke-fitting cast around Troy \newline
 the elder Atreides ruler, fortunate man \\[1mm] % %
 τοιόνδε Τροίαι περιβαλὼν ζευκτήριον \newline
 ἄναξ Ἀτρείδης πρέσβυς εὐδαίμων ἀνὴρ.\end{quote} %
} %

The Latin loanword ἀρτέμων / \textit{artemōn} / “foresail” used in Acts %
apparently crossed into literary Greek usage contemporary or subsequent to the reign of Aurelius. %
Given the Athenian dramatic sensibilities elsewhere, %
and Atreidean subtexts throughout, %
its adjacent term ζευκτήριος / \textit{dzeukērios} / “yoke-fit” %
likely also reflects a learned allusion to the Aeschylean term. %
\clearpage

\noindent{}\textbf{ἀρχιποίμην} / \textit{archipoimēn} / “chief-shepherd” appears just once in the NT, %
in 1Pet 5:4. %
Its first occurrence in non-Christian literature belongs to Herodianʼs authentic \textit{General Prosody}: %

\vspace*{4mm} 
 \begin{quote} Τὰ εἰς ην ἀπὸ ῥητῶν ἤτοι διαλεγομένων τῶν εἰς ην βαρύνεται %
 οἷον Ἕλλην φιλέλλην, Πανέλληνες ὡς Παναχαιοί ἀπὸ μέρους τὸ %
 πᾶν ὡς παρὰ τῷ ποιητῇ «Λέλεγας καὶ Τρῶας ἀναίρειν» (\textit{Il}. 20, 96).
 Ἕλληνες γὰρ παρ’ αὐτῷ οἱ τὴν Φθιῶτιν Ἑλλάδα κατοικοῦντες %
 «Μυρμιδόνες δὲ καλεῦνται καὶ Ἕλληνες καὶ Ἀχαιοί» (\textit{Il}. 2, 684), %
 εἴρην μελλείρην, ποιμήν \textbf{ἀρχιποίμην}, αὐχήν ἐριαύχην, %
 ῥήν πολύρρην παρὰ τῷ τὰ Ναυπακτικὰ ποιήσαντι· (\textit{Pros.cath.} \textit{GG} 3.1:16) \\[2mm] %
 Words ending with ην, whether based on standard terms or discursive ones %
 that end with ην, are accented, such as Hellene philhellene, %
 Panhellenes as Panachaeans, as the whole derives from the part, %
 as with the poet, “To destroy Leleges and Troy”. %
 For with him the Hellenes are those who inhabit Phthiotin Hellas, %
 “Myrmidons they are called, and Hellenes and Achaeans”, %
 [Lacedaemonian] young-man near-man, shepherd \textbf{chief-shepherd}, neck arch-necked, %
 sheep [\textit{rhēn}] Polyrrhenian, with the one who created the \textit{Naupatica}.\end{quote} % 
 \vspace*{4mm}

The passage is unclear as to whether or not Herodian attests the word ἀρχιποίμην / \textit{archipoimēn} / “chief-shepherd” %
for the largely lost text of the \textit{Naupatica} (6th BCE). %
Thus we are left with Herodian himself as the first clear, non-Christian witness to the lemma. %
It also appears in the \textit{Testament of Judah}. %
Like the other books in the \textit{Testaments of the Twelve Patriachs}, %
this one is viewed by experts as a Christian composition of the second half of the 2nd century CE.\endnote{%
De Jonge 1985:85.} %
\clearpage

 \begin{quote}Ἦν δέ μοι καὶ κτήνη πολλά, %
 καὶ εἶχον \textbf{ἀρχιποίμενα} Ἵραν τὸν Ὀδολαμίτην· (\textit{Test.Jud.} 4.8.1) \\[0.4mm] %
 Many herds belonged to me, 
 and I had Ira the Adullamite as \textbf{chief-shepherd}. \end{quote} % 
 \vspace*{3mm}

\noindent{}This midrash expands on LXX Gen 38:1--20, %
where Judah has intercourse with his daughter-in-law, Tamar, %
and arranges for a goat to be sent to her by way of the %
Adullamite previously identified as “Ira” (38:1). %

In I Peter, the word occurs within a set of moral guidance for leaders in the Christian community. %
This passage (1Pet 5:2--4) can be found as quoted in the chapter on Philostratus II (222--249 CE) in regard to the \textit{Heroikos}. %
For two apparent neologisms, one evidenced by Herodian and the other by Philostratus, %
to appear in close proximity strengthens the case for I Peter being redacted quite late, %
sometime between the late 2nd century and the second quarter of the 3rd century. %
While the latter option seems quite late, %
it is right in the neighborhood of Origen of Alexandria and his commentary on Luke, %
which is the earliest Christian text to evidence the possible reception of 1Pet 5:2--4.\endnote{%
Preserved in a catena fragment as \textit{Luc.cat.} fr. 58a--c (GCS 49:251): %
Ποιμέσι πρῶτον ἀποκαλύπτεται τὸ ὑπὸ τῶν ἀγγέλων ὑμνούμενον
μυστήριον οἵτινες τύπον ἐπέχουσι τῶν τὰς ἐκκλησίας μελλόντων ποιμαίνειν.
καὶ μάλα εἰκότως. (58b) αὐτοὺς γὰρ ἔδει τό· ‹Ἐπὶ γῆς εἰρήνη› κατὰ τὴν πρώτην τάξιν ἀκούειν,
οἵτινες ἤμελλον τὴν εἰρήνην ἐπιφωνεῖν παντὶ τῷ τῆς ἐκκλησίας πληρώματι.
τὸ δὲ τῶν ποιμένων πρόσωπον καὶ ἡ δι’ ἀποκαλύψεως αὐτοῖς γενομένη
χαρὰ σημαίνει σαφῶς, ὡς ἐπὶ ‹τὸ πλανώμενον πρόβατον› ἦλθεν ‹ὁ ποιμὴν
ὁ καλός›. (58c) ποιμένας γὰρ οὐδὲν οὕτως εὐφραίνει, ὡς ἡ τοῦ ἀπολωλότος βοσκήματος
εὕρεσις, ὅπερ οὐκ ἦν ἑτέρου τινὸς εὑρεῖν ἢ ‹τοῦ ἀρχιποιμένος Χριστοῦ›.} %
It also happens to align with the mid-3rd century CE dates for three different Greek papyri %
that contain the lemma.\endnote{%
sb 8 9908 (267 CE); sb 8 9909 (267 CE); p.prag. 2 126 (166--269CE), written by a secondary hand at the later end of the date range.} \\[1.5mm] %

\noindent{}\textbf{ἐφευρετής} / \textit{epheuretēs} / “inventor” %
has 326 hits in the TLG. %
Its first potential non-Christian occurrence is within %
the famous collection of Greek poems known as the \textit{Anacreonta}, %
specifically in fragment 38, an incomplete drinking song to Bacchus.\endnote{%
West 1984, fr. 38, line 3: %
\vspace*{1mm} 
 \begin{quote}Ἱλαροὶ πίωμεν οἶνον, \newline
 ἀναμέλψομεν δὲ Βάκχον, \newline
 τὸν ἐφευρετὰν χορείας, \newline
 τὸν ὅλας ποθοῦντα μολπάς, \newline
 τὸν ὁμότροπον Ἐρώτων, \newline
 τὸν ἐρώμενον Κυθήρης· \\[1mm] %
 Cheerful let us drink wine, \newline
 we shall raise a strain to Bacchus, \newline
 the inventor of choral dance, \newline
 he who craves all songs, \newline
 the same manner of Lusts, \newline
 which is lusted by Cytherea. %
 \end{quote} % 
 \vspace*{1mm}
}
This collection, however, is certainly not by Anacreon himself, %
but rather clearly Hellenistic and eclectic, even though “few” of its %
poems “give any clue to the date of their composition”.\endnote{Campbell 1988:10.} %
Based on reception and citation, %
there was some sort of \textit{Anacreonta} collection of poems circulating in the mid- to late 2nd century CE, %
a collection known to Gellius, Clement of Alexandria, and Hippolytus of Rome. %
But what they attest for that collection bears no match to what remains in todayʼs primary witness, %
the Palatine manuscript.\endnote{Campbell 1988:12--13.} %
Furthermore, in a study of the vocabulary of the Palatine \textit{Anacreonta}, %
Brioso Sánchez recounted this specific lemma (ἐφευρετής / \textit{epheuretēs}) as a sign of late compositional work, %
since the lemma is not attested before the 1st century CE.\endnote{Campbell 1988:14, n18; Brioso Sánchez 1970.} %
Other terms in Palatine \textit{Anacreonta} appear nowhere else before the 2nd, 3rd, and even 4th century CE (14). %
Based on a careful study of vocabulary and syntax, Brioso Sánchez dated fr. 38 to the range 400--600 CE.\endnote{%
Campbell 1988:17; Brioso Sánchez 1970.} %
Thus the \textit{Anacreonta} TLG result is invalidated as a viable neologistic anchor. %
\vfill{}
The lemma did not turn up in searches of major databases of digitized Greek papyri or Greek inscriptions, %
but it does occur in two works attributed to Herodian, both likely spurious on the whole. %
First we will examine the passage from Παρεκβολαὶ τοῦ μεγάλου ῥήματος / %
\textit{Excursus de Verbo Magno} / \textit{Excerpts on the Great Verb}. %
While Dickey catalogs this work under the spuria of Herodian,\endnote{Dickey 2014:330.} %
this only reflects the general characterization made by the editor, La Roche, %
who noted in the introduction that the style is uncharacteristic of Herodian, %
and thus that the work is generally spurious, likely compiled by a Byzantine grammarian. %
At the same time, La Roche adds that this epitomizer pulled extensively from Apollonius and Herodian, %
yet also that any modern effort to extract authentic grains from the chaff inevitably bears little fruit.\endnote{La Roche 1863:3.} %
\vfill{}
Caveats and uncertainties notwithstanding, %
the relevant quotation belongs to a section on “Definite Passives” / Ὁριστικὰ παθητικά, %
where [Ps-]Herodian explores words derived from the stem for the verb “to find” / εὑρίσκω / \textit{heuriskō}: %

\clearpage
 \begin{quote}Πόθεν τὸ εὑρέθην; ἀπὸ τοῦ εὑρῶ εὑρήσω εὕρηκα, %
 καὶ συστολῇ τοῦ η εἰς τὸ ε εὑρέθην, %
 ὅθεν τὸ ῥηματικὸν ὄνομα εὑρετὴς καὶ \textbf{ἐφευρετής}.\endnote{Παρεκβολαὶ τοῦ μεγάλου ῥήματος in La Roche 1863:24.} \\[1mm] %
 Whence [comes] “I was found” (\textit{heurethēn})? %
 From “I find” (\textit{heurō}), “I will find” (\textit{heurēsō}), “I have found” (\textit{heurēka}), %
 and by the contraction of the η into ε, “I was found” (\textit{heurethēn}), %
 whence [comes] the verbal noun “finder” (\textit{heuretēs}) and “\textbf{inventor}” (\textit{\textbf{epheuretēs}}). \end{quote} % 
 \vspace*{3mm}

The other ostensibly Herodian attestation also belongs to a spurious work: %
Περὶ παραγωγῶν γενικῶν ἀπὸ διαλέκτων / \textit{De Derivationibus Genitivis ex Dialectis} / %
\textit{On Genitive Derivations by Dialect}.\endnote{Dickey 2014:338.} %

\vspace*{3mm} 
 \begin{quote}Πότε τὰ εἰς σ λήγοντα φυλάττουσιν ἐν τῇ συνθέσει τὴν τάσιν τοῦ ἁπλοῦ; %
 ἡνίκα εὑρεθῶσιν ὑπὲρ μίαν συλλαβήν· %
 οἷον, ἐρανιστὴς, αὐχερανιστής· %
 εὑρετὴς, \textbf{ἐφευρετής}· \\[1mm] %
 When do word endings with \textit{sigma} in the compound keep %
 the pitch of the singular? %
 Whenever they are found with more than one syllable, %
 such as: feast-club-member (\textit{eranistēs}), high-feast-club-member (\textit{aucheranistēs}); %
 finder (\textit{heuretēs}), \textbf{inventor} (\textit{\textbf{epheuretēs}}). \end{quote} % 
 \vspace*{3mm}

The lemma also appears in Alexander Aphrodisiasʼ late 2nd, early 3rd century CE \textit{Commentary on Aristotleʼs Metaphysics}, %
but apparently in a spuriously attributed portion, at least according to the TLG work breakdown.\endnote{%
Under the heading Book N of Aristotle’s \textit{Metaphysics} in Hayduck 1891:794. %
} %
In the mid-3rd century, it is used several times by Origen of Alexandria, %
and also by Porphyry in his \textit{Contra Christianos} (fr. 15). %

The lemma does not seem to appear among digitized Greek papyri or Greek inscriptions. %
Given the vague or problematic evidence, %
a precise point of origin is not attainable. %
Even so, it seems reasonable to suggest that this lemma first emerged in the mid- to late-2nd %
century, if not the early 3rd century. %

Its solitary appearance in the NT is in canonical Romans 1:30, %
in a portion likely missing from Marcionʼs \textit{Apostolos}. %

\vspace*{2mm} 
 \begin{quote} παρέδωκεν αὐτοὺς ὁ θεὸς εἰς ἀδόκιμον νοῦν, ποιεῖν τὰ μὴ καθήκοντα, %
 ²⁹ πεπληρωμένους πάσῃ ἀδικίᾳ πονηρίᾳ πλεονεξίᾳ κακίᾳ, %
 μεστοὺς φθόνου φόνου ἔριδος δόλου κακοηθείας, % 
 ³⁰ ψιθυριστάς, καταλάλους, θεοστυγεῖς, ὑβριστάς, ὑπερηφάνους, %
 ἀλαζόνας, \textbf{ἐφευρετὰς} κακῶν, γονεῦσιν ἀπειθεῖς \\[1mm] %
 God handed them over to a valueless mind, to do improper things, %
 ²⁹ having been filled with all injustice, evil, malice, avarice, %
 full of ill-will, murder, strife, baiting, evil-disposition, %
 ³⁰ whisperers, slanderers, God-scorns, insolents, super-proud, %
 imposters, \textbf{inventors} of evils, disobedient to parents. \end{quote} % 
 \vspace*{2mm}

 This same section has two other potential neologisms noted in the penultimate chapter \textit{infra}: %
 ἀντιμισθία (Rom 1:27, also in Theophilus of Antioch) and κατάλαλος (Rom 1:30, also in \textit{Hermas}). %
 Perhaps the vice list reflected the spontaneous creativity of the editor of canonical Romans. %
 Even so, a mid- to late-2nd century origin best fits the evidence. \\[1mm] %

\noindent{}\textbf{ζιζάνιον} / \textit{dzidzanion} / “darnel”, %
according to the LSJ entry, refers to \textit{Lolium temulentum}, %
a wheat-like grain that inebriates (i.e., poisons) its consumers. %
Its ancient roots go back to the Sumerian word for wheat, \textit{zizân}. %
The Greek term may have its first %
attestation in the long-form \textit{Partitiones} of [Ps-]Herodian.\endnote{%
Boissonade 1819:40.}
\vspace*{3mm}

 \begin{quote}Τὰ ἀπὸ τῆς ζι συλλαβῆς ἀρχόμενα διὰ τοῦ ι γράφεται· %
   οἷον· \textbf{ζιζάνιον}, ἡ ἄκανθα· ζιβύνιον, ἀκόντιον ὁλοσίδηρον· καὶ τὰ λοιπά. \clearpage %
  [Words] that begin with the syllable \textit{zi} are written with an \textit{iota}, %
  such as: \textbf{\textit{zizanion}}, or weed; \textit{zibunion}, a wholly-iron javelin; and the rest.\end{quote} % 
\vspace*{3mm} 

A confounding example appears in the \textit{Apocalypse of Moses}, %
a text in Latin known as the \textit{Life of Adam and Eve}. %
Here Satan says to the serpent in the garden: %

\vspace*{3mm} 
 \begin{quote}διὰ τί ἐσθίεις ἐκ τῶν \textbf{ζιζανίων} τοῦ Ἀδάμ %
 καὶ οὐχὶ ἐκ τοῦ παραδείσου; (\textit{ApMos} 16.5)\endnote{%
 Text from Tromp 2005:138.} \\[1mm] %
 “For what reason do you eat from the \textbf{\textit{zizania}} of Adam, %
 and not from the [fruit] of paradise?”\end{quote} % 
 \vspace*{3mm} 

\noindent{}The date of the Greek text of \textit{ApMos} is highly uncertain and speculative. %
If it was originally a Hebrew composition, as Johnson speculated, then the Hebrew text %
may have dated between 100 BCE and 200 CE, “but more probably toward the end of the first Christian century”, %
yet “the Greek and Latin texts” date “between that time and A.D. 400”.\endnote{%
Johnson 2005:252.} 

Within the canonical NT, %
the term ζιζάνιον appears eight times, %
all of them within a canonical Matthean fable (13:25--27, 29--30) and %
subsequent fable explanation (13:36, 38, 40). %
This fable lacks an overall parallel in Marcionʼs *Ev or the other three NT Gospels. %
The Matthean examples use the term specifically in relation to an “enemy” who sows %
\textit{zizania} in the wheat fields of his rival, %
and the subsequent effort by the field-owner to deal with this problem. % 
The explanation ties the story to final judgment and spiritual actors (God, the evil one, angels). %

It appears in other Christian writings sometimes considered to be early, %
i.e., the 13-letter collection of the letters of Ignatius of Antioch, %
specifically in Ign13 \textit{Ephesians} 11.9.1, in a portion missing from Syriac Ign3 and Greek Ign7,\endnote{%
For the Syriac three letter collection at this point, see Cureton 1849:23--26. %
For the Greek seven letter collection at this point, see Ehrman 2003:230--231.} %
and thus likely to date well after the mid-2nd century. %
A Greek fragment (fr. 28) of the fourth book of Irenaeusʼ \textit{Against Heresies} %
that quotes the Matthean fable provides possible confirmation %
of the late 2nd century (ca. 180 CE) Christian usage of the term. %
The compact clustering of the unusual term in Matthew, %
its relatively late reception, %
and the earliest possible non-Christian appearance of the lemma in the Aurelian period, %
all point to the Matthean usage as a late-Aurelian or post-Aurelian insertion. %
Metatextually, it may comment on textual contaminations coming up alongside %
the emerging Christian sacred scriptures themselves, %
and/or an admission that divinely-oriented patience with heretical enemies, %
rather than human efforts to control the editing of scriptures, %
was now the only viable path. \\[1.5mm] %

\noindent{}\textbf{θρῆσκος}|\textbf{θρησκός} / \textit{thrēskos} / “religious” %
only appears 34 times in the TLG. %
Here again [Ps-]Herodian in the long-form \textit{Partitiones} may provide %
the earliest attestation to a potential neologism:\endnote{%
Boissonade 1819:59. %
The cognate term θρησκεία appears in several Greek papyri dating from %
41 CE (p.lond. 6 1912) through the 6th century, %
as well as in a Greek inscription later falsely ascribed to Nazareth, %
NAZA0002 (27 BCE -- 100 CE, slab, imperial edict, %
on which see \url{search.inscriptionsisraelpalestine.org/inscriptions/naza0002})} %

\vspace*{3mm} 
 \begin{quote}Πλὴν τοῦ θρῆνος, ὁ κλαυθμός· θρηνῳδῶ· θρηνῳδία· %
 θρηνῳδός· θρῆνυς, τὸ ὑποπόδιον· Θρήϊκες, οἱ Θρᾷκες· θρησκεύω· % 
 θρησκεία· \textbf{θρῆσκος}, ὁ περί τι πιστός· καὶ Θρῆσσα, ἡ Θρᾳκική. \\[1mm] %
 In addition there is lamentation, wailing; to sing dirges; dirge-song; %
 dirge-singer; the \textit{threnus}, a sandal; Thracians, those of Thracia; %
 to be religious; religion; \textbf{religious}, one faithful to something; and Thracianne, %
 what is Thrace-like. \end{quote} % 
 \vspace*{3mm} 

The same rare term only appears once in the canonical NT, %
specifically in the pseudonymous letter of James (1:26): %

\vspace*{3mm} 
 \begin{quote}\textls[-20]{Εἴ τις δοκεῖ \textbf{θρησκὸς} εἶναι, %
  μὴ χαλιναγωγῶν γλῶσσαν αὐτοῦ ἀλλὰ ἀπατῶν καρδίαν αὐτοῦ, %
  τούτου μάταιος ἡ θρησκεία. %
  ²⁷ θρησκεία καθαρὰ καὶ ἀμίαντος παρὰ τῷ θεῷ καὶ πατρὶ αὕτη ἐστίν, %
  ἐπισκέπτεσθαι ὀρφανοὺς καὶ χήρας ἐν τῇ θλίψει αὐτῶν, %
  ἄσπιλον ἑαυτὸν τηρεῖν ἀπὸ τοῦ κόσμου.} \\[1mm] %
 If anyone thinks to be \textbf{religious} % 
 while not bridling his tongue, but deceiving his heart, %
 the religion of such a one is empty. %
 ²⁷ Religion clean and spotless with the God and Father is this, %
 to look after orphans and widows in their affliction, %
 to keep oneself stainless from the world. %
 \end{quote} % 
 \vspace*{3mm}

\noindent{}The usage of the apparent neologism is interesting, %
but far more striking is the closely connected, %
parallel use of a second cognate word, θρησκεία / “religion”, %
just as [Ps-]Herodian used at this precise point in a lexicographical compendium. %
Indeed, read James 1:27 once more. %
It is structured like a dictionary entry, is it not? %
Given that the long-form \textit{Partitiones} derives from around the %
time of Marcus Aurelius, %
the most reasonable date for this well-integrated usage in James, %
together with this entire moral treatise-as-letter, %
is Aurelian or post-Aurelian. \\[2mm] %

\noindent{}\textbf{νομοδιδάσκαλος} / \textit{nomodidaskalos} / “law-teacher” first appears in TLG results %
for non-Christian texts in the long-form \textit{Partitiones} of [Ps-]Herodian.\endnote{%
Boissonade 1819:92.} %

\vspace*{3mm} 
 \begin{quote}Πᾶσα λέξις ἀπὸ τῆς νο συλλαβῆς ἀρχομένη διὰ τοῦ %
 ο μικροῦ γράφεται· οἷον· νόμος· νομικός· \textbf{νομοδιδάσκαλος}· %
 νομίζω· νόμισμα· νόσος· νοσῶ· νόσημα· νόται, %
 παρὰ Ῥωμαίοις τὰ γράμματα, ἐξ οὗ καὶ νοτάριος, ὁ γραμματικός· \\[1mm] %
 Every word that begins with the syllable \textit{vo} is written with the short \textit{ο}. %
 For example: law; legal; \textbf{law-teacher}; to reckon; currency; %
 disease; to be sick; illness; “letters” (\textit{notai}) among Romans, %
 from which also comes notary, the letter-skilled. \end{quote} % 
 \clearpage

 \noindent{}The same lemma does not seem to appear in any Greek papyri or Greek inscriptions, %
 but it does appear three times in the NT. First in our list is Luke 5:17, %
a verse most likely missing from Marcionʼs \textit{Evangelion}:

\vspace*{3mm} 
 \begin{quote}Καὶ ἐγένετο ἐν μιᾷ τῶν ἡμερῶν καὶ αὐτὸς ἦν διδάσκων, %
   καὶ ἦσαν καθήμενοι Φαρισαῖοι καὶ \textbf{νομοδιδάσκαλοι} %
   οἳ ἦσαν ἐληλυθότες ἐκ πάσης κώμης τῆς Γαλιλαίας %
   καὶ Ἰουδαίας καὶ Ἰερουσαλήμ· \\[1mm] %
   And it happened on one of the days when he was also teaching, %
   and there were seated Pharisees and \textbf{law-teachers} %
   who had arrived from every village of the Galilee %
   and Judea and Jerusalem. %
   \end{quote} % 
   \vspace*{3mm}
  
\noindent{}Sounds like a Tannaitic council of Rabbis, %
and similar to the Jerusalem Council described in Acts 15, doesnʼt it? % %
No surprise, the next occurrence of the lemma is in Acts (Acts 5:34): %

\vspace*{3mm}
 \begin{quote}ἀναστὰς δέ τις ἐν τῷ συνεδρίῳ Φαρισαῖος
 ὀνόματι Γαμαλιήλ, \textbf{νομοδιδάσκαλος} τίμιος παντὶ τῷ
 λαῷ \\[1mm] %
 Then someone, after rising up in the Sanhedrin, %
 a Pharisee, Gamaliel by name, a \textbf{law-teacher} esteemed by all the people. \end{quote} % 
 \vspace*{3mm}

\noindent{}Finally, the same lemma appears in 1Tim 1:6--7:

\vspace*{3mm} 
 \begin{quote} ὧν τινες ἀστοχήσαντες ἐξετράπησαν εἰς ματαιολογίαν %
  ⁷ θέλοντες εἶναι \textbf{νομοδιδάσκαλοι}, %
  μὴ νοοῦντες μήτε ἃ λέγουσιν μήτε περὶ τίνων διαβεβαιοῦνται. \\[1mm] %
  Some of them, becoming misaligned, went astray into empty-speech, %
  ⁷ wishing to be \textbf{law-teachers}, %
  not understanding either what they are saying or about what they are thoroughly assured. \end{quote} % 
 \clearpage

 Previous generations of scholars have found points of commonality between %
 the Pastorals and Acts, or more broadly between the Pastorals and Luke-Acts. %
 Our finding for this lemma fits into those arguments, %
 while pushing the respective compositional and/or editorial of all of them to around %
 or after the time of Marcus Aurelius. %
 It should finally be noted that a TLG search of the lemma νομοδιδάσκαλος %
 does indeed turn up one hit prior to Herodian, %
 specifically in Plutarch (\textit{Cato Maior} 20.6), but the stem is distinctive: %
νομοδιδάκτης / \textit{nomodidaktēs} / “law-teacher”, %
not νομοδιδάσκαλος / \textit{nomodidaskalos} / “law-teacher”.\endnote{%
The late 2nd c. Artemidorus of Daldis also uses νομοδιδάκτης (\textit{Oneir}. 2.29).} \\[1mm] %

\noindent{}\textbf{σθενόω} / \textit{sthenoō} / “to strengthen” %
appears just once in the NT, as the closing benediction of I Peter 5:10. %

\vspace*{2mm} 
 \begin{quote} Ὁ δὲ θεὸς πάσης χάριτος, ὁ καλέσας ὑμᾶς %
 εἰς τὴν αἰώνιον αὐτοῦ δόξαν ἐν Χριστῷ, %
 ὀλίγον παθόντας αὐτὸς καταρτίσει, στηρίξει, \textbf{σθενώσει}, θεμελιώσει.
 ¹¹ αὐτῷ τὸ κράτος εἰς τοὺς αἰῶνας· ἀμήν. \\[1mm] %
 Now the God of all favor, who called you into his eternal glory in Christ, %
 after you have suffered a little, %
 will Himself supply, establish, \textbf{strengthen}, foundation-set [you]. %
 ¹¹ To Him [be] the power into the ages. Amen. \end{quote} % 
 \vspace*{2mm}

\noindent{}The lemma does not occur again in a Christian text until Gregory Nazianzenʼs %
 \textit{Orations} 42 (\textit{Supremum vale}), %
 which clearly refers back to the benediction in I Peter. %
 It also does not appear in the major databases of Greek papyri or Greek inscriptions.\endnote{%
 Note that IG II\textsuperscript{2} 4514 — IG II\textsuperscript{3},4 837, Attica, Athens, mid-2nd century, %
 hymn to Asclepius, clearly attests σθένω / \textit{sthenō} / “to have strength”, %
 not σθενόω / \textit{sthenoō} / “to strengthen”.}

The first hit in a non-Christian text is in the authentic \textit{General Prosody} of Herodian, %
in which he clearly distinguishes between the two cognate and similarly spelled lemmas: \clearpage
 
 \begin{quote}Τὰ εἰς νω παραληγόμενα τῷ ε, περισπᾶται, εἰ προκατάρχοι ὄνομα,
 ἐξ οὗ γέγονε \newline
 φρενός φρενῶ, \newline %
 ἑνός ἑνῶ, \newline %
 ξένος ξενῶ, \newline %
 ἀσθενής ἀσθενῶ, \newline %
 σθένος \textbf{σθενῶ}, \newline %
 στένος στενῶ. \\[1mm]%
\textls[15]{τὸ γὰρ στένω οὐκ ἔχει ὄνομα οὔτε μένω, %
 μεθ’ ὧν καὶ φένω τὸ φονεύω καὶ πένω, %
 ἀφ’ οὗ πένομαι, καὶ σθένω καὶ θένω τὸ τύπτω. (\textit{Pros.cath.}, \textit{GG} 3.1:449)}\\[2mm] %
 Those that end with νω (\textit{nō}) that are preceded with the ε (\textit{e}) %
 are circumflexed if they should first cause a noun. %
 From this came \\[1mm]
 from “mind” (\textit{phrenos}), “to be mindful” (\textit{phrenō}); \newline
 from “one” (\textit{henos}), “to unify” (\textit{henō}); \newline
 from “stranger” (\textit{xenos}), “to be a stranger” (\textit{xenō}); \newline
 from “weak” (\textit{asthenēs}), “to be weak” (\textit{asthenō}); \newline
 from “strength” (\textit{sthenos}), “\textbf{to strengthen}” (\textit{\textbf{sthenō}}); \newline
 from “narrow” (\textit{stenos}), “to narrow” (\textit{stenō}). \\[1.5mm]
 \textls[-15]{The [verb] “to narrow” (\textit{stenō}) lacks a noun, %
 as does “to remain” (\textit{menō}). %
 Along with these “to slay” (\textit{phenō}) came “to murder” (\textit{phoneuō}) %
 and “to toil” (\textit{penō}), from which came “to be poor” (\textit{penomai}), %
 and [from] “to be strong” (\textit{sthenō}) and “to smite” (\textit{thenō}) [came] “to strike” (\textit{tuptō}).} %
 \end{quote} %
 \vfill{}

\noindent{}Given the rarity of the lemma, the precision of Herodianʼs explanation, %
and his fame and authority as a grammarian in Rome, %
Herodian is more likely a progenitor and/or popularizer of this neologism, %
and I Peter is more likely secondary and subsequent. \clearpage

\noindent{}\textbf{σιμικίνθιον} / \textit{simikinthion} / “apron” %
is a loanword from the Latin \textit{semicinctium}. %
It appears just once in the NT, in Acts 19:12. %
It provides a bit of dramatic costumery %
and may reflect the emergence of Christian martyr relics %
akin to the sort described in the \textit{Martyrdom of Polycarp}. %
In its literary context, the term makes vivid the miraculous power of Paul, whose skin %
merely had to come into contact with ordinary pieces of cloth %
in order for divine healing to be conveyed through them. %

\vspace*{3mm} 
 \begin{quote}Δυνάμεις τε οὐ τὰς τυχούσας ὁ θεὸς ἐποίει διὰ τῶν χειρῶν Παύλου, %
 ¹² ὥστε καὶ ἐπὶ τοὺς ἀσθενοῦντας ἀποφέρεσθαι ἀπὸ τοῦ χρωτὸς %
 αὐτοῦ σουδάρια ἢ \textbf{σιμικίνθια} καὶ ἀπαλλάσσεσθαι ἀπ’ αὐτῶν τὰς νόσους, %
 τά τε πνεύματα τὰ πονηρὰ ἐκπορεύεσθαι. \\[1mm] %
 Powers, and not those by happenstance, God did through the hands of Paul, %
 ¹² so that even when kerchiefs or \textbf{aprons} were carried off from his skin %
 [and placed] upon those who were sick, %
 it made the diseases release from them, %
 and also made the evil spirits leave.  \end{quote} % 
 \vspace*{3mm} 
 
 It it not until the mid-4th century (Cyril of Jerusalem, \textit{Cat.illum.} 10.19) %
 that the lemma re-emerges in Christian literature in the TLG. %
 In non-Christian Greek literature in the TLG, %
 this term only occurs in two works, %
 the [post-]Aurelian long-form \textit{Partitiones} of [Ps-]Herodian, %
 and the 5th--6th century lexicon of Hesychius. %
Hesychius has several additional lemma entries preceding and following σιμικίνθιον / \textit{simikinthion} %
that are not evident in the passage of [Ps-]Herodian. %
By my lights, though there may be some later word forms and Byzantine glosses mixed among the parallels and definitions, %
there is no clear indication of Christian interpolation. %
Quite the opposite, actually, since there are multiple indications %
of reverence for traditional Greek and Roman divinities. \clearpage

The line breaks in the following excerpt are ours, %
meant to facilitate correlation between the Greek \textit{Partitiones} and English translation. %
\vspace*{2mm}

 \begin{quote}σιγῶ· σιγή· σιωπή· \newline %
 σίδηρος· σιδηρόφρακτος, ὁ ὡπλισμένος· \newline %%
 Σιδὼν, πόλις, καὶ κλίνεται Σιδῶνος· Σίδης, ἡ πόλις· \newline %%
 Σιβύλλα, κύριον; σίελος, τὸ σάλος· \newline %
 σίκερα, πόμα σκευαστόν· σικινὶς, εἶδος ὀρχήσεως· \newline %%
 σίκυα, τὰ τετράγγουρα, ὅθεν καὶ σικυήλατον· \newline %%
 σίμβλον, τὸ τῶν μελισσῶν οἴκημα· \newline %%
 \textbf{σιμικίνθιον}, τὸ φακιόλιον· \newline %%
 σιμὸς, ὁ πατζός· Σιμόεις, ποταμός· \newline %%
 σίνω, τὸ βλάπτω· \newline %%
 σίνηπι· Σινώπη, χώρα· Σινωπεὺς ἀνήρ· \newline %
 σιπύη, ἡ ἀρτοθήκη· \newline %
 Σιχὰρ καὶ Σισάρα, κύρια· καὶ Σισυφός· (Boissonade 1819:124--126) \\[1mm]
 To be silent; Silence; Stillness; \newline %
 Iron; Iron-clad, one who is armored; \newline %
 Sidon, a city, and it is inflected \textit{Sidōnos}; Sides, the city; \newline %
 Sibyl, divinity; saliva, the rolling swell; \newline %
 Liquor, a crafted drink; Sikinnis, a form of dancing; \newline %
 Cucumbers, the large cucumbers, whence also cucumber-bed; \newline %
 Beehive, the dwelling of the bees; \newline %
 \textbf{Apron}, the face-cloth; \newline
 Simian, the snub-nosed; Simoeis, a river; \newline
 To hurt, that is, to damage; \newline %
 Mustard; Sinope, a region; Sinopean, a man [from Sinope]; \newline %
 A meal-tub, a bread-container; \newline %
 Sichar and Sisera, divinities; and Sisyphus.%
 \end{quote} % 
 \clearpage

Could the lemma in this context reflect the existence and textual contents of a lexicon %
in use in Rome around the time of Marcus Aurelius? %
A few of the terms may be of later provenance, %
but most were well known by that time. 
It is difficult to be certain, %
but it seems reasonable that σιμικίνθιον / \textit{simikinthion} / “apron” appeared %
in a lexicon in the 170s CE, or not long after. %
As we stated before, if the \textit{Partitiones} reflect later interpolations, %
that tends to strengthen the case for a later date for the corresponding NT usage of that word. %

What seems quite reasonable is that Acts is unlikely %
the earliest extant text to attest this Latin word being imported into Greek usage. %
It is more likely such a word would be preserved in a Greek grammatical work such as this, %
one that reflects the enormous breadth of vocabulary in use among Greek-speaking elites. %
To put a finer point on it, while direct dependence for this word is far from conclusive, %
it is more likely that Acts was indebted to an introductory/exhaustive word-list similar to that of [Ps-]Herodian %
than for this rare Latin loanword to have been inserted as a gloss into this text. \\[1mm] %%  

\noindent{}\textbf{σουδάριον} / \textit{soudarion} / “towel|napkin” %
derives from the Latin \textit{sudarium}. %
It has 279 hits at present in the TLG, %
but none of these securely precede the 2nd century CE. %
The lemma occurs twice in the authentic \textit{General Prosody} of Herodian (\textit{GG} 3.1:365, 534), %
and is also found in the authentic Περὶ διχρόνων / \textit{On Anceps Vowels} (\textit{GG} 3.2:13). %
On the latter work, Dyck and Dickey concur that its content was excerpted from the \textit{General Prosody} %
and is thus authentic to Herodian. %
Other non-Christian authors around or after Herodianʼs work %
also attest the lemma, including Julius Pollux (\textit{Onomast}. 7.71) and Aelius Promotus (\textit{Potency} 21.1). %
Pollux specifically distinguishes previous terms for a particular type %
of garment from \textit{soudarion} as a contemporary term: \clearpage %

 \begin{quote}τὸ δὲ ἡμιτύβιον, ἔστι μὲν καὶ τοῦτο Αἰγύπτιον, %
 εἴη δ’ ἂν κατὰ τὸ ἐν τῇ μέσῃ κωμῳδίᾳ καψιδρώτιον καλούμενον,
 ὃ νῦν \textbf{σουδάριον} ὀνομάζεται\endnote{Bethe 1900.2:72, Citing III 466.325 Ko. to confirm the Middle Comedy term.} \\[1mm] %
 But the “linen cloth” (\textit{hēmitubion}), while this is also Egyptian; %
 but it might be what is in Middle Comedy called “napkin” (\textit{kapsidrōtion}), %
 which now is called \textit{\textbf{soudarion}}.\end{quote} % 
\vspace*{2mm}

In Greek papyri and inscriptions, there are several to consider, %
but apparently none before 138 CE, or even before Herodian.\endnote{%
LSJ lists Corpus Papyrorum Raineri [CPR] 1.27.7 (ii A.D.), which is stud.pal.20.15 — SPP 20 15 — chrest.mitt.289 — TM 15012 — Oct 27, 189 CE, %
or “AD 190 Arsinoite” in chrest.mitt.289. %
LSJ also lists Supp.Epigr.7.417 (Dura), which in SEG online is SEG 50-1394 — Dura-Europas. The ‘archive’ of the merchant Aurelius Nebouchelos, 235--240 A.D., %
and in PHI is PH320658 — SEG 7:417 — Syria-Middle Euphrates — Dura-Europos (Salhiye, nr.) — 235--240 AD — Excav. Dura-Europos 4 (1933) 98, 227. %
Another Arsinoite/Fayum fragment is also from 230 CE: SPP 20 31 — CPR 1 124 — Stud. Pal. 20 31 — TM 15019.}
The lemma appears twice in the canonical NT, %
once in Luke 19:20, in the so-called Parable of the Pounds, %
but specifically in a verse unattested for and likely missing from Marcionʼs \textit{Evangelion}. %
It also appears in John 11:44, in the story of the resurrection of Lazarus. \\[1mm] %

\noindent{}\textbf{συκομορέα} / \textit{sukomorea} / “sycamore-tree” %
has 95 hits in the TLG. %
In Christian tradition, the lemma is famously associated with the dimunutive Zacchaeus %
who climbed this type of tree to get a look at Jesus in canonical Luke 19:4. %
While that episode is attested for Marcionʼs \textit{Evangelion}, %
this particular lemma is not. %
In non-Christian texts, it first appears in the authentic work of Herodian \textit{On Orthography}, %
within the context of an alphabetized list clarifying penultimate vowels: %

\vspace*{1mm} 
 \begin{quote}συβόσιον: ι. \textbf{συκομορέα}: ε. συκοφαντία: ι (\textit{Orthog}. GG 3.2:585)\\[1mm] %
 herd of swine (\textit{subosion}): \textit{iota}. \textbf{Sycamore-tree} %
 (\textit{\textbf{sukomorea}}): \textit{epsilon}. Brown-nosing (\textit{sukophantia}): \textit{iota}.\end{quote} % 
 \vspace*{1mm}
  %
\noindent{}\textls[-15]{The same lemma is also in [Ps-]Herodian, \textit{Partitiones} (Boissonade 1819:217).} %

In Christian authors, it first appears in Origen, %
who uses the term twice in reference to the canonical Lukan tale (\textit{Ps.hom.} 24.5), %
connecting it to the LXX term συκάμινος / \textit{sukaminos} / “mulberry-tree”. %
LSJ notes this traditional and more common spelling, as well as συκόμορος / \textit{sukomoros}, %
the latter evident in the 1st century CE in Strabo and Dioscorides Pedanius, %
and as far back as the 3rd century BCE PCair.Zen. 83.3. %
Among manuscripts of canonical Luke, there seems to be quite a bit of confusion in spelling, %
particularly whether to use an \textit{omega} or \textit{omicron} between the \textit{mu} and \textit{rho}, %
and also whether the penultimate vowel should be an \textit{epsilon} or \textit{alpha-iota} diphthong.\endnote{%
Swanson 1995:320.} Yet all these variants do not register either of the two traditional forms noted in LSJ. %
That is to say, canonical Luke almost certainly evidences the new spelling first documented securely by Herodian, %
around the time of Marcus Aurelius, %
lending further support to a post-Aurelian date for a major round of editing to canonical Luke. \\[1mm] %

\noindent{}\textbf{ὑπερεκτείνω} / \textit{huperekteinô} / “to stretch out beyond measure” %
has 101 TLG hits, but no hits in searches of major databases of Greek papyri or inscriptions. %
Its first non-Christian occurrence is Herodianʼs authentic Περὶ Ὀδυσσειακῆς προσῳδίας / \textit{On the Prosody of the Odyssey}.\endnote{%
On its authenticity, see Dickey 2014:337 and Dyck 1993:786.} %
In context, the excerpt attempts to explain an unusual word from the opening of \textit{Odyssey} book 23, %
specifically \textit{Od} 23.3, where Eurycleia rushes (as much as an old woman can rush) %
to tell Penelope that Odysseus has returned. %

\vspace*{2mm} 
 \begin{quote}ὑπερικταίνοντο: Ἀρίσταρχος μὲν ἄγαν ἐπάλλοντο καὶ ἐκινοῦντο %
 προθυμουμένης αὐτῆς βαδίζειν ταχέως, μὴ δυναμένης δέ, ἀλλὰ κατὰ %
 βραχὺ διὰ τὸ γῆρας. οἱ δὲ \textbf{ὑπερεξετείνοντο}. \\[1mm] %
 “were hobbling along” (\textit{hyperiktainonto}): Aristarchus\endnote{Aristarchus of Samothrace.} %
 interpreted it as “they were quivering” (\textit{epallonto}) and “they were moving” (\textit{ekinounto}), %
 as she [the old woman] was eager to proceed quickly, but was not capable, %
 except little by little due to old age. %
 Others [read] “\textbf{were stretching out beyond measure}”. \end{quote} % 
 \clearpage

Herodian goes on to give his preferred interpretation. %
The anonymous attribution, though, %
should give us pause to consider whether the word was newly minted by Herodian, %
or whether it represented one of his contemporary grammarians or perhaps even his grammarian father. %
In any case, the same unusual, and perhaps new word used by Herodian appears %
in canonical 2Cor 10:14, in an extended portion absent from Marcionʼs \textit{Apostolos}. %

\vspace*{2mm} 
 \begin{quote}οὐ γὰρ ὡς μὴ ἐφικνούμενοι εἰς ὑμᾶς \textbf{ὑπερεκτείνομεν} ἑαυτούς, %
 ἄχρι γὰρ καὶ ὑμῶν ἐφθάσαμεν ἐν τῷ εὐαγγελίῳ τοῦ Χριστοῦ· \\[1mm] %
 \textls[-15]{For we do not as [if] not reaching to you \textbf{extend ourselves beyond measure}, %
 for as far as you we also outran in the evangelion of the Christ.}\end{quote} % 
 \vspace*{2mm}

\noindent{}Against the background of Herodianʼs comment, %
the Homeric allusion becomes crystal clear. %
Canonical Paul (now implicitly picturing himself as elderly, together with his team), %
takes on the persona of Herodianʼs Eurycleia %
and (perhaps satirically) turns the episode into a metaphor of athletic competition. %

This new formulation likely helps explain another cognate neologism found just below in canonical 2Cor 10:16, %
\textbf{ὑπερέκεινα} / \textit{huperekeina} / “beyond measure”, %
which we have listed under Clement of Alexandria in the penultimate chapter. %
While Clement may well have coined the latter in his editing of canonical 2Cor, %
he very likely owed his use of the lemma ὑπερεκτείνω / \textit{huperekteinô} to Herodian. \\[1mm] %

\noindent{}\textbf{χοϊκός} / \textit{choikos} / “dust-born” | “pertaining to dust” appears first in extant texts either in %
1Cor 15:47--49 (Marcionite and canonical), [Ps-]Herodian, \textit{Partitiones}, or a Greek inscription: IG II\textsuperscript{2} 13139/42 — PH15633 — Attic — post. mid-2nd century CE. %
By around 180 CE in Irenaeusʼ major work, it was popular and core to emerging proto-Orthodox %
theology of divine incarnation (\textit{AH} 1.10--11, 14, 16). %
The lemma also appears in a relatively late (probably ca. 3rd century CE), Christian portion of \textit{SibOrac} 8, at line 460, %
with an incarnational theology that resonates with that of Irenaeus.\endnote{The early history of this lemma %
requires more thorough research than can be covered here in the first edition.} %